% Generated by report/tools/tables.py from the result files of the repository. Do not edit.
\begin{table}[tbp]
\centering
\footnotesize
\setlength{\tabcolsep}{4pt}
\caption[Les chiffres de Pineau, comme contexte]{Les chiffres de Pineau (2018), comme contexte et jamais comme verdict : ses parts sur 50 articles de RL de 2018, pour les six critères qu'elle chiffre (pourcentage publié, \file{study/results\_final.json} ; entre parenthèses le $k/n$ des tests, son pourcentage multiplié par 50 et arrondi, soit 56~\% pour Q5 et 4~\% pour Q8), comparées à notre lecture de 50 articles de 2018 avec notre grille, puis aux 45 articles expérimentaux du premier tirage de 2026, grands modèles de langage compris. La première ligne de chaque critère mesure l'écart entre notre grille et la sienne sur la même année : notre lecture est plus sévère sur Q3 et plus large sur Q7, ce qui a conduit à ne fonder les verdicts que sur la comparaison avec la même grille (écart 20 de \file{study/PROTOCOL.md}, tableau~\ref{tab:sondage-pineau}). La seconde mêle le changement des pratiques, celui de la grille et celui de la population. Écart en points avec l'intervalle de Newcombe à 95~\%, test exact de Fisher, correction de Holm sur Q3 à Q8 dans chaque famille de comparaisons. \textbf{En gras}, ce qui passe la correction de Holm. Le tirage de Pineau reste inconnu. Sources : \file{study/extension\_stats.json}, \file{study/results\_final.json}.}
\label{tab:sondage-pineau-contexte}
\begin{tabularx}{\linewidth}{@{}l>{\raggedright\arraybackslash}Xrrr r@{\hspace{0.28em}}l rr@{}}
\toprule
Code & Critère, puis comparaison & $k/n$ & \% & Pineau 2018 & \multicolumn{2}{c}{Écart [Newcombe]} & $p$ & $p$ Holm \\
\midrule
Q3 & \multicolumn{8}{@{}l}{Tous les hyperparamètres de la méthode} \\
 & \quad notre lecture de 2018 & 26/47 & 55 & 90 (45/50) & \textbf{−35} & \textbf{[−50\,;\,−17]} & \textbf{0,0002} & \textbf{0,001} \\
 & \quad premier tirage 2026 & 30/45 & 67 &  & \textbf{−23} & \textbf{[−39\,;\,−7]} & \textbf{0,006} & \textbf{0,02} \\
\addlinespace
Q4 & \multicolumn{8}{@{}l}{Tous les hyperparamètres des références} \\
 & \quad notre lecture de 2018 & 17/48 & 35 & 60 (30/50) & −25 & [−42\,;\,−5] & 0,02 & 0,06 \\
 & \quad premier tirage 2026 & 8/44 & 18 &  & \textbf{−42} & \textbf{[−57\,;\,−22]} & \textbf{5\,×\,10\textsuperscript{−5}} & \textbf{0,0002} \\
\addlinespace
Q5 & \multicolumn{8}{@{}l}{Code en lien} \\
 & \quad notre lecture de 2018 & 15/48 & 31 & 55 (28/50) & −25 & [−42\,;\,−5] & 0,02 & 0,06 \\
 & \quad premier tirage 2026 & 30/45 & 67 &  & +11 & [−9\,;\,+29] & 0,30 & 0,91 \\
\addlinespace
Q6 & \multicolumn{8}{@{}l}{Choix des hyperparamètres décrit} \\
 & \quad notre lecture de 2018 & 12/47 & 26 & 20 (10/50) & +6 & [−11\,;\,+22] & 0,63 & 1 \\
 & \quad premier tirage 2026 & 12/45 & 27 &  & +7 & [−10\,;\,+23] & 0,47 & 0,95 \\
\addlinespace
Q7 & \multicolumn{8}{@{}l}{Évaluation sur conditions tenues à part} \\
 & \quad notre lecture de 2018 & 16/48 & 33 & 10 (5/50) & \textbf{+23} & \textbf{[+7\,;\,+39]} & \textbf{0,006} & \textbf{0,03} \\
 & \quad premier tirage 2026 & 31/44 & 70 &  & \textbf{+60} & \textbf{[+42\,;\,+73]} & \textbf{1\,×\,10\textsuperscript{−9}} & \textbf{6\,×\,10\textsuperscript{−9}} \\
\addlinespace
Q8 & \multicolumn{8}{@{}l}{Test statistique appliqué} \\
 & \quad notre lecture de 2018 & 3/48 & 6 & 5 (2/50) & +2 & [−8\,;\,+13] & 0,67 & 1 \\
 & \quad premier tirage 2026 & 2/45 & 4 &  & 0 & [−10\,;\,+11] & 1 & 1 \\
\bottomrule
\end{tabularx}
\end{table}
